\section{Những phép đo chưa ai chạy, và giá của chúng}
\label{sec:ch18-phep-do-chua-ai-chay}

\index{khoảng trống nghiên cứu!phép đo chưa ai chạy|(}%
Năm điều kiện vừa xếp ra loại bớt chứ không cộng thành một thiết kế. Mục này đi
từ ràng buộc sang việc: trên đường đi, sách gặp bốn phép đo mà không ai chạy, và
mục này xếp chúng theo giá chứ không theo mức thú vị.

\index{khoảng trống nghiên cứu!phép đo rẻ nhất là đọc lại một bảng đã in}%
Phép đo rẻ nhất không đòi chạy gì cả, và
mục~\ref{sec:ch18-ba-bai-le-ra-da-dong} vừa đi qua nó. Bảng của bài chuẩn đã in
sẵn, cùng với một cột ngẫu nhiên đặt làm mốc, và thư viện sinh ra nó là thư viện
công khai. Việc chưa ai làm là đọc bảng ấy theo cột thay vì theo hàng: với mỗi
chỉ số, hỏi nó có tách được cột mốc ngẫu nhiên khỏi các phương pháp đã công bố
hay không, và tách được bao xa. Đọc theo hàng thì được câu hỏi phương pháp nào
tốt hơn; đọc theo cột thì được câu hỏi chỉ số nào phân biệt được cái gì. Cùng một
bảng trả lời được cả hai câu, và bài chỉ hỏi câu thứ nhất.

Phép đo ấy rẻ tới mức đáng ngờ, nên phải nói ngay nó không kết luận được gì. Một
chỉ số không tách được mốc ngẫu nhiên khỏi các phương pháp trên \emph{một} mô
hình, \emph{một} bộ dữ liệu tổng hợp và \emph{một} mức tương quan đặc trưng thì
mới chỉ hỏng ở đó. Đó là một điều kiện cần theo đúng nghĩa mà
mục~\ref{sec:ch08-dieu-kien-can} đặt ra, và một điều kiện cần bị vi phạm thì bác
được, còn được thỏa thì không xác nhận gì. Cái nó cho là một câu hỏi chạy được
ngay với chi phí gần bằng không.

\index{khoảng trống nghiên cứu!phép đo thứ hai, nút thắt khái niệm}%
Phép đo thứ hai đắt hơn một chút và
mục~\ref{sec:ch15-doc-duoc-dieu-khien-duoc} đã tính giá của nó: dựng một tập khái
niệm không cần nhãn thì làm được, còn hỏi mỗi đơn vị bên trong mô hình có ép được
đầu ra hay không thì tốn đúng một lượt chạy cho mỗi đơn vị. Chưa ai chạy nó ở quy
mô mà câu trả lời có nghĩa.

\index{khoảng trống nghiên cứu!phép đo thứ ba, trên mô hình thật}%
Phép đo thứ ba là phép đo mà mục~\ref{sec:ch14-dieu-chua-giai-quyet} nêu, và giá
của nó là giá của việc dựng một thiết lập. Cả ba bài của
chương~\ref{ch:giai-thich-duoi-tan-cong} làm việc trên những thiết lập mà ground
truth đã biết vì có người dựng ra nó, dù người ấy là bên tấn công hay là người
viết phương trình sinh dữ liệu. Không bài nào hỏi một mô hình bình thường, huấn
luyện trên dữ liệu thật, thì phần sai lệch của lời giải thích nằm ở đâu và lớn
tới đâu. Đó cũng là điều kiện thứ nhất của
bảng~\ref{tab:ch18-nam-dieu-kien} đọc thành một thí nghiệm.

\index{khoảng trống nghiên cứu!phép đo thứ tư, cái đắt nhất}%
Phép đo thứ tư là phép đo đắt nhất, và nó là phép đo mà cả chương này tồn tại để
chỉ vào. Nó là phép đo của bài neo, chạy ở cột thứ ba của
hình~\ref{fig:ch18-hinh-dang-khoang-trong}: dựng một bộ nhiệm vụ mà đầu ra để lộ
phép tính đã sinh ra nó, gán nhãn theo thiết kế ấy, rồi chấm từng chỉ số của
chương~\ref{ch:do-do-trung-thuc} xem điểm của nó có tách được lời giải thích đúng
cơ chế khỏi lời giải thích không đúng cơ chế hay không. Giá của nó là công dựng
bộ nhiệm vụ, và mục~\ref{sec:ch16-gia-duong-dung-san} đã đo phạm vi mà một bộ như
vậy với tới trên chính bài neo, nên cái giá ấy đã biết trước chứ không phải đoán.

\index{khoảng trống nghiên cứu!một phép đo đã được chạy và trả lời}%
Trong danh sách này có đúng một mục đã được gạch đi, và gạch nó là việc của
chương~\ref{ch:con-nguoi-vang-mat}. Mục~\ref{sec:ch09-doi-tuong-khac-nhau} kể ba
thứ đi qua được ranh giới, và thứ thứ ba là một phép kiểm mà phía attribution
chạy được ngay: đo mức đồng thuận giữa các chỉ số, việc không cần ground truth
nào. Nó đã được chạy, và
mục~\ref{sec:ch12-ket-qua-va-ranh-gioi} thuật lại kết quả qua bài báo dẫn nó: các
chỉ số độ trung thực của attribution đặc trưng chỉ tương quan yếu với nhau và
\emph{có thể} cho ra những bảng xếp hạng phương pháp mâu thuẫn nhau. Chỗ rào ấy
là chỗ rào của bài được dẫn chứ không của sách, và nó phải giữ nguyên. Câu trả
lời đi cùng chiều với câu trả lời ở phía chuỗi suy luận.

Chỗ ấy phải đọc cho đúng mức, vì nó dễ bị đọc quá tay. Các chỉ số bất đồng với
nhau thì chưa nói chỉ số nào sai, và
mục~\ref{sec:ch09-ket-qua} đã ghi hai cách giải thích cùng khớp với mức đồng
thuận thấp: chúng đo những tính chất khác nhau, hoặc chúng cùng nhắm một đại
lượng với nhiễu lớn. Điều phép kiểm ấy làm được là loại bỏ khả năng thứ ba, rằng
các chỉ số đo cùng một thứ và đo đúng. Nó thu hẹp chứ không kết luận.

\index{khoảng trống nghiên cứu!phép đo chưa ai chạy|)}%
Đọc bốn phép đo theo giá thì thấy một thứ tự ngược với thứ tự mà lĩnh vực đi. Cái
rẻ nhất chưa ai làm, dù số liệu đã in ra từ năm 2021. Cái đắt nhất cũng chưa ai
làm, và nó là cái duy nhất trong bốn cái trả lời thẳng câu hỏi của
định nghĩa~\ref{def:ch01-do-trung-thuc}. Cái đã được làm là cái không đòi ground
truth, và nó cho một câu trả lời thu hẹp chứ không kết luận. Nói cách khác, lĩnh
vực đã làm đúng phần việc không đòi trả giá, và ba phần còn lại đều đang chờ ở
chỗ chúng được đặt xuống.
